Setjum $O_1$ og $O_2$ sem miðju hringjanna. Ath. að $\triangle O_1AE$ og $\triangle O_1CE$, og $\triangle O_2BE$ og $\triangle O_2DE$ eru eins. Þá er $EA=EC$ og $EB=ED$.
\newline Fáum: $EA+\frac12EA=6$, $EB+2EB=6$ og $CD=EC-ED$ þ.a. $CD=4-2=2$.

\begin{center}
\begin{tikzpicture}[
    dot/.style={circle, fill=black!80, draw=black, line width=0.3pt, inner sep=1pt},
    lbl/.style={font=\small, inner sep=2pt}]
  % --- data: centres and radii (change these freely) ---
  \def\xa{0}    \def\ya{0}     \def\ra{1.30}  % small circle
  \def\xb{3.55} \def\yb{-0.28} \def\rb{2.08}  % large circle

  % --- tangent directions ---
  \pgfmathsetmacro{\dist}{veclen(\xb-\xa,\yb-\ya)}
  \pgfmathsetmacro{\ang}{atan2(\yb-\ya,\xb-\xa)}
  \pgfmathsetmacro{\tout}{\ang+acos((\ra-\rb)/\dist)} % external tangent
  \pgfmathsetmacro{\tin}{\ang-acos((\ra+\rb)/\dist)}  % internal tangent

  \coordinate (O1) at (\xa,\ya);
  \coordinate (O2) at (\xb,\yb);
  \coordinate (A) at ($(O1)+(\tout:\ra)$);
  \coordinate (B) at ($(O2)+(\tout:\rb)$);
  \coordinate (C) at ($(O1)+(\tin:\ra)$);
  \coordinate (D) at ($(O2)+(\tin+180:\rb)$);
  \coordinate (E) at (intersection of A--B and C--D);

  % --- drawing ---
  \draw[black!45] (O1) circle (\ra);
  \draw[black!45] (O2) circle (\rb);
  \draw[black!70, thick] 
    (A) -- (B) 
    (C) -- (E)
    (A) -- (O1) -- (C)
    (B) -- (O2) -- (D);
  \draw[black, dotted, thick]
    (O1) -- (E) -- (O2);

  \node[dot] at (A) {}; \node[lbl, above] at (A) {$A$};
  \node[dot] at (B) {}; \node[lbl, above] at (B) {$B$};
  \node[dot] at (E) {}; \node[lbl, above] at (E) {$E$};
  \node[dot] at (D) {}; \node[lbl, left]  at (D) {$D$};
  \node[dot] at (C) {}; \node[lbl, below]  at (C) {$C$};
  \node[dot] at (O1) {}; \node[lbl, left]  at (O1) {$O_1$};
  \node[dot] at (O2) {}; \node[lbl, below left]  at (O2) {$O_2$};
\end{tikzpicture}
\end{center}
